\cotefile{1}
\scanpage{79}

\uline{Corollaire 4} Soit $u$ un op. de Fredholm d'un espace de Hilbert $E$. Alors

\[
\boxed{|\alpha_n(u)|\leq\alpha_n(\deleted{\unsure{\sqrt{u^*u}}}|u|\deleted{{}_{\unsure{2}}})}\tag*{(18)}
\]

(\deleted{En} d'autres termes) les coefficients du développement de Fredholm $\det(1+zu)$ sont majorés en valeur, terme à terme, par les coefficients du déterminant de Fredholm $\det(1+z\unsure{h|H|\lvert u\rvert})$ relatif à l'hermitien positif $h=\sqrt{\unsure{uu^*|u^*u}}$. On en conclut en effet

\uline{Corollaire 5} Dans les conditions précédentes, on a

\[
|\det(1+zu)|\leq\det(1+|z|h)\quad\deleted{\unsure{(\sqrt{u^*u})}}\tag*{(19)}
\]

\unsure{Prouvons|Montrons} (18). On a \inserted{en effet}, $|\alpha_nu|=|\mathrm{Tr}_n\,\hat u|\leq\|\hat u_n\|_1=\unsure{\mathcal U_n|\mathcal A_n|\alpha_n}(\sqrt{u^*u})$ (formule (16)).

\deleted{\uline{Théorème 2} Soit $u$ une \deleted{application}\inserted{op.} de Fredholm dans l'espace de Hilbert $H$. Alors la suite \inserted{$(\lambda_i)$} des valeurs propres de $u$ est sommable \deleted{\unsure{et}} \deleted{\unsure{(1)}} \deleted{\illegible{4}} \deleted{\unsure{valeurs de $u$}}, on a}

\[
\boxed{\deleted{\det(1+u)=\prod_i(1+\lambda_i)}}\tag*{\deleted{(20)}}
\]

\deleted{\uline{Corollaire 1} Le déterminant de Fredholm $\det(1+zu)$ de $u$ est une fonction entière de genre $0$\,;}

\[
\deleted{\det(1+zu)=\prod_i(1+z\lambda_i)}\tag*{\deleted{(22)}}
\]

\deleted{\inserted{Appliquant (21) à $zu$ au lieu de $u$, on trouve}}

% margin note written sideways
\margin{\deleted{Le corollaire 1 \unsure{ci-}\unsure{dessus} équivaut \deleted{\unsure{aux}}\inserted{à la} formule, (23) $\alpha_n(u)=\sum_{i_1<\cdots<i_n}\lambda_{i_1}\deleted{\unsure{(u)}}\ldots\lambda_{i_n}\deleted{\unsure{(u)}}$ dont le \unsure{plus intéressant|premier cas particulier} est (23) $\boxed{\mathrm{Tr}\,u=\sum\lambda_i(u)}$}}

\deleted{\uline{Démonstration} du th. 2. \deleted{\unsure{Nous prouvons}} \unsure{On prouve} d'abord le lemme (20) \unsure{sur les valeurs}, \unsure{nous bornant à} la formule (20) : soit un \unsure{op. et un autre}, \illegible{3} \unsure{propre}, soit \unsure{une projection} \illegible{2}, alors tout \unsure{revient} qu'il en est \unsure{ainsi}, et que \illegible{1} est \unsure{identique} dès que \deleted{\unsure{tout}} \unsure{la} \illegible{1} est \unsure{nulle}.}

% archive sheet number: 80
